%# -*- coding:utf-8 -*-
\documentclass[10pt,aspectratio=169,mathserif]{beamer}		
%设置为 Beamer 文档类型，设置字体为 10pt，长宽比为16:9，数学字体为 serif 风格

%%%%-----导入宏包-----%%%%
\usepackage{zju}			%导入 zju 模板宏包
\usepackage{ctex}			%导入 ctex 宏包，添加中文支持
\usepackage{amsmath,amsfonts,amssymb,bm}   %导入数学公式所需宏包
\usepackage{color}			 %字体颜色支持
\usepackage{graphicx,hyperref,url}
\usepackage{metalogo}	% 非必须
%% 上文引用的包可按实际情况自行增删
%%%%%%%%%%%%%%%%%%
\usepackage{fontspec}
\usepackage{xeCJK}
% \setCJKmainfont{Source Han Sans SC}

\beamertemplateballitem		%设置 Beamer 主题

%%%%------------------------%%%%%
\catcode`\。=\active         %或者=13
\newcommand{。}{．}				
%将正文中的“。”号转换为“.”。中文标点国家规范建议科技文献中的句号用圆点替代
%%%%%%%%%%%%%%%%%%%%%

%%%%----首页信息设置----%%%%
\title[浙江大学 Beamer 模板]{浙江大学 Beamer 模板}
\subtitle{——这里是副标题}			
%%%%----标题设置

\author[Guang Touge]{
  光头哥 \\\medskip
  {\small \url{chenqiyuan1012@foxmail.com}} \\
  {\small \url{https://www.zju.edu.cn/}}}
%%%%----个人信息设置
  
\institute[IOPP]{
  计算机学院 \\ 
  浙江大学}
%%%%----机构信息

\date[Aug. 31 2023]{
  2023年8月31日}
%%%%----日期信息
  
\begin{document}

\begin{frame}
	\titlepage
\end{frame}				%生成标题页

\section{提纲}
\begin{frame}
	\frametitle{提纲}
	\tableofcontents
\end{frame}				%生成提纲页

%%%%========== 第一部分：问题 ==========%%%%
\section{问题}
\begin{frame}
	\frametitle{问题}
	\begin{itemize}
		\item 研究背景
		\item 核心问题是什么？
		\item 为什么这个问题重要？
	\end{itemize}
\end{frame}

%%%%========== 第二部分：相关工作 ==========%%%%
\section{相关工作}
\begin{frame}
	\frametitle{相关工作}
	\begin{itemize}
		\item 已有方法一
		\item 已有方法二
		\item 已有方法的局限性
	\end{itemize}
\end{frame}

%%%%========== 第三部分：挑战/创新 ==========%%%%
\section{挑战/创新}
\begin{frame}
	\frametitle{挑战与创新}
	\begin{itemize}
		\item 现有方法面临的挑战
		\item 本文的创新点
		\item 主要贡献
	\end{itemize}
\end{frame}

%%%%========== 第四部分：方法 ==========%%%%
\section{方法}
\begin{frame}
	\frametitle{方法}
	\begin{itemize}
		\item 整体框架
		\item 关键模块/算法
		\item 理论分析
	\end{itemize}
\end{frame}

%%%%========== 第五部分：实验 ==========%%%%
\section{实验}
\begin{frame}
	\frametitle{实验}
	\begin{itemize}
		\item 实验设置（数据集、基线方法、评估指标）
		\item 主要实验结果
		\item 消融实验
	\end{itemize}
\end{frame}

%%%%========== 第六部分：思考 ==========%%%%
\section{思考}
\begin{frame}
	\frametitle{思考}
	\begin{itemize}
		\item 论文的优点与不足
		\item 可能的改进方向
		\item 对未来研究的启发
	\end{itemize}
\end{frame}

\end{document}